\chapter{Simple Linear Regression}

\section{Predictor and Response}\index{simple linear regression}\index{regression!simple linear}
\newterm{Simple linear regression} summarizes the relationship between two
quantitative variables. The \index{predictor variable}\newterm{predictor} $x$
(also called the explanatory or independent variable) is used to predict the
\index{response variable}\newterm{response} $y$ (also called the outcome or
dependent variable). It is \emph{simple} because there is only one predictor.
Multiple regression, covered later, uses two or more predictors.

\section{The Model}
A straight line is fit to the scatterplot:
$$\hat{y} = b_0 + b_1 x$$
Here $b_0$ is the intercept and $b_1$ is the slope, which is the estimated
change in $y$ for a one-unit increase in $x$. The population model adds an error term:
$$y = \beta_0 + \beta_1 x + \epsilon$$
The error $\epsilon$ is the part of $y$ that the line does not capture.

\section{Least Squares}\index{least squares}
The line is chosen to minimize the sum of the squared vertical distances from the
points to the line. Those vertical distances are the \index{residual}\newterm{residuals}
$e_i = y_i - \hat{y}_i$. The formulas that minimize $\sum e_i^2$ are
$$b_1 = \frac{S_{xy}}{S_{xx}}, \qquad b_0 = \bar{y} - b_1 \bar{x}$$
where $S_{xx} = \sum (x_i - \bar{x})^2$ and $S_{xy} = \sum (x_i - \bar{x})(y_i - \bar{y})$.

\subsection{A Worked Example}
Five students studied for $x = 1, 2, 3, 4, 5$ hours and scored
$y = 2.1, 3.9, 6.2, 7.8, 10.1$ (out of 10 points on a quiz).
We have $\bar{x} = 3$ and $\bar{y} = 6.02$. Also $S_{xx} = 10$ and
$S_{xy} = 19.9$. So
$$b_1 = \frac{19.9}{10} = 1.99, \qquad b_0 = 6.02 - 1.99(3) = 0.05$$
The fitted line is $\hat{y} = 0.05 + 1.99x$. For a student who studies 6 hours,
the prediction is $\hat{y} = 0.05 + 1.99(6) = 11.99$. Be careful: this is a
prediction for $x = 6$, which is just outside the data. Predictions far outside
the range of the data are unreliable.

The residuals are $0.06, -0.13, 0.18, -0.21, 0.10$. They always sum to zero for
a least squares line with an intercept, because the line passes through
$(\bar{x}, \bar{y})$.

\section{Interpreting the Slope}
The slope $b_1 = 1.99$ means that each additional hour of study is associated
with about 1.99 more points on the quiz. The word \emph{associated} matters. A
regression line describes a pattern in the data. It does not prove that studying
causes the score to rise, unless the data came from a well-designed experiment.

\begin{Exercise}[label=slr1, title=Predict a value]
Using $\hat{y} = 0.05 + 1.99x$, predict the score for a student who studies
4 hours.
\end{Exercise}
\begin{Answer}[ref=slr1]
$\hat{y} = 0.05 + 1.99(4) = 8.01$.
\end{Answer}

\begin{Exercise}[label=slr2, title=Find a residual]
For the student who studied 3 hours, the observed score was 6.2. What is the residual?
\end{Exercise}
\begin{Answer}[ref=slr2]
The predicted score is $0.05 + 1.99(3) = 6.02$. The residual is $6.2 - 6.02 = 0.18$.
\end{Answer}

\begin{Exercise}[label=slr3, title=Interpret the intercept]
Explain why the intercept $b_0 = 0.05$ is not useful for interpretation in this example.
\end{Exercise}
\begin{Answer}[ref=slr3]
The intercept is the predicted score for 0 hours of study. No data was collected
near $x = 0$, and a student who studies zero hours could still score something
other than 0.05, so the intercept is an extrapolation and is not a meaningful
estimate by itself.
\end{Answer}
